% 第八章：结论与改进路线

\section{结论与改进路线} \label{sec:conclusion}

\subsection{结论}

\paragraph{实现了什么。}一个内置于 Linux 6.8 的调度类：策略号 8，2752 行 C，配 29 个对外声明、3 个 tracepoint、28 项 \texttt{debugfs} 计数器与 7 个可独立应用的补丁。它不是原型片段——能编进 \texttt{bzImage}、能在 QEMU 里起来、能让任务真正在这个策略下跑完并交回正确的计数。

\paragraph{它确实在工作。}512 次唤醒触发 487 次选核（其余走 wakelist 快路径），271 次完整决策全部落地，接管率 $97.9\%$，平均可运行核数 4.0。近一半唤醒（44.4\%）由"目标核还空着"直接复用，既省决策又保局部性。这组数字的价值在于它排除了两类常见假阳性和假阴性：回调被亲和性锁死（\texttt{shared} 配置下 \texttt{select\_calls}=0 就是反面证据），以及"算出来但没落地"。

\paragraph{代价还没收敛。}主口径下 Yat 的墙钟是 CFS 的 \textbf{2.00} 倍（5 次重复的中位数），当前实现不具备可用性。归因不是猜的：pin 到单核（回调一次都不被调用）仍慢 1.70 倍，说明主体代价来自"唤醒时进入这个调度类"，选核只占边际（$2.00/1.70\approx1.18$）；没有唤醒时两者没有可分辨差异（1.10，跨 1），说明代价随唤醒次数增长；\texttt{balance\_runs}=757 $>$ 唤醒 512 次，说明几乎每个调度点都在重扫就绪池。

\paragraph{修掉了一处真实的并发缺陷。}锁反序由 \texttt{lockdep} 实证、由 \texttt{task\_dead} 的 \texttt{irqsave} 修复，修复后同类告警计数归零，且没有引入性能回归。这一条是本文里唯一"改完就能验证对错"的结论，因此也最硬。

\paragraph{推翻了项目早期的一组结论。}前期用 Java 模拟器得到的"11.66\% Makespan 改进、30.21 个百分点命中率提升、96\% 胜率"不可复现，根因是两套模型口径、基线状态泄漏与评分函数缺项三层叠加。仿真器已从证据链中移除，只作为设计期校验留在仓库文档里。这件事本身就是工程结论：\textbf{一个只在理想模型里被验证过的调度改进，不足以支撑"在真实系统上有效"的结论}。

\subsection{改进路线}

以下四条按"改动范围—预期收益—风险"排序，可独立采纳。

\paragraph{路线 A：把唤醒路径的重复扫描去掉（改动小，收益最大）}
\texttt{balance\_runs}=757 而唤醒只有 512 次，说明几乎每个调度点都在扫描全部 DAG 的就绪池。两步：给就绪池加一层"按 DAG 索引"的间接表，把全表扫描降为按 DAG 定位；把 balance 的触发从"每个调度点"改为"被唤醒任务的亲和集与上次不同"时才做。验收标准是计数器而非墙钟：\texttt{balance\_runs} 应降到与唤醒同量级，主口径中位数先逼近 \texttt{shared} 的 1.70，再压到 1.3 以下。这一条不需要改算法，只需要改数据结构。

\paragraph{路线 B：把决策的候选集收窄（改动中，收益可测）}
当前 MSF 在全部可运行核上计算 benefit，容易把任务推离刚放过同 DAG 的核。可引入"负载窗口"：候选集限定为 $\ell_k\le\ell_{\min}+\alpha\hat{C}$，在窗口内再按 benefit 选。这与前期仿真里验证过的思路同源，但仿真不能作为证据，必须在核心里用 \texttt{place\_hit}/\texttt{miss} 与迁移计数重新验证。$\alpha$ 需要扫描，且要防止退化成纯 WFD。

\paragraph{路线 C：离开 TCG（改动小，结论价值大）}
所有实测都在 QEMU/TCG 上，墙钟与真实机器不可比，且同配置波动 28\%。在裸机上重跑同一套脚本（\texttt{--accel kvm} 或物理机），才能回答"慢一倍"是不是模拟器的产物。这是把本文从"能跑"推进到"可发表"的必要一步。

\paragraph{路线 D：标定 CRP 的归一化窗口（改动小，影响深远）}
\texttt{yat\_get\_crp\_ratio()} 的 \texttt{max\_recency} 实际是 1\,ms，而本轮负载的重用距离在 $2\times10^5$\,ns 量级，尚未饱和。一旦睡眠间隔变长或核数增加，CRP 项会静默退化为 0，MSF 变成纯负载均衡——这种失效没有任何报错。标定方法是：在 $10\,\mu$s--$10$\,ms 之间扫睡眠间隔，记录 $\rho$ 的分布与饱和点，再定常量。在没有标定数据之前不应改这个数值。

\subsection{验收清单}

交付验收见\autoref{tab:accept}。

\begin{table}[H]
\centering
\caption{交付验收}
\label{tab:accept}
\small
\begin{tabular}{@{}p{5.4cm}p{8.4cm}@{}}
\toprule
项 & 状态 \\
\midrule
\texttt{make verify-all} & 通过（静态门禁 + 符号一致性 + 补丁包 + 文档链接） \\
调度类可编入 \texttt{bzImage} 并在 QEMU 启动 & 通过（4 vCPU，\texttt{nr\_cpu\_ids}=4） \\
回调被调用 & 通过（\texttt{select\_calls}=487/512） \\
决策落地 & 通过（\texttt{jobs\_placed}=271，接管率 97.9\%） \\
锁反序 & 已修复，告警 0 次，补丁 0007 可独立应用 \\
代价可复现 & 三配置 $\times$ 5 次，中位数与全距见第\ref{sec:test}节 \\
可用性（性能） & \textbf{未达标}：主口径约 2.00$\times$ CFS \\
真实机器验证 & \textbf{未做}（路线 C） \\
\bottomrule
\end{tabular}
\end{table}